%%%%%%%%%%%%%%%%%%%%%%%%%%%%%%%%%%%%%%%%%%%%%%%%%%%%%%%%%%%%%%%%%%%%%%
% Lớp 8 | Chương X. Một số hình khối trong thực tiễn | Hình chóp tam giác đều
% Dạng 8H0?38-2: Diện tích xung quanh, thể tích hình chóp tam giác đều
%%%%%%%%%%%%%%%%%%%%%%%%%%%%%%%%%%%%%%%%%%%%%%%%%%%%%%%%%%%%%%%%%%%%%%


% nguồn: De_AI_Toan8_cuoi_hoc_ki_i_1 (nhập 21/11/2026)
\begin{bt}%[8H0N38-2]%[ma:d306b0df]
	Cho hình chóp tam giác đều có độ dài cạnh đáy bằng $6\text{ cm}$ và độ dài trung đoạn bằng $8\text{ cm}$.
a) Tính nửa chu vi đáy của hình chóp.
b) Tính diện tích xung quanh của hình chóp tam giác đều đó.
	\loigiai{
		a) Chu vi đáy của hình chóp là $C = 3 \cdot 6 = 18\text{ cm}$. \\
Nửa chu vi đáy là $p = \dfrac{18}{2} = 9\text{ cm}$. \\
b) Diện tích xung quanh của hình chóp tam giác đều là: \\
$S_{xq} = p \cdot d = 9 \cdot 8 = 72\text{ cm}^2$.
	}
\end{bt}

% nguồn: De_AI_Toan8_cuoi_hoc_ki_i_1 (nhập 21/11/2026)
\begin{bt}%[8H0N38-2]%[ma:60353683]
	Tính diện tích xung quanh của một hình chóp tam giác đều biết độ dài cạnh đáy bằng $10\text{ cm}$ và độ dài trung đoạn bằng $12\text{ cm}$.
	\loigiai{
		Nửa chu vi đáy của hình chóp tam giác đều là: \\
$p = \dfrac{3 \cdot 10}{2} = 15\text{ cm}$. \\
Diện tích xung quanh của hình chóp là: \\
$S_{xq} = p \cdot d = 15 \cdot 12 = 180\text{ cm}^2$.
	}
\end{bt}

% nguồn: De_AI_Toan8_cuoi_hoc_ki_i_1 (nhập 21/11/2026)
\begin{bt}%[8H0N38-2]%[ma:dd22743a]
	Tính thể tích của một hình chóp tam giác đều có diện tích mặt đáy bằng $24\text{ cm}^2$ và chiều cao bằng $9\text{ cm}$.
	\loigiai{
		Áp dụng công thức tính thể tích hình chóp tam giác đều $V = \dfrac{1}{3} S_{\text{đáy}} \cdot h$, ta có: \\
$V = \dfrac{1}{3} \cdot 24 \cdot 9 = 72\text{ cm}^3$. \\
Vậy thể tích của hình chóp là $72\text{ cm}^3$.
	}
\end{bt}

% nguồn: De_AI_Toan8_cuoi_hoc_ki_i_1 (nhập 21/11/2026)
\begin{bt}%[8H0N38-2]%[ma:43de8795]
	Tính thể tích của một hình chóp tam giác đều có diện tích mặt đáy bằng $36\text{ cm}^2$ và chiều cao bằng $10\text{ cm}$.
	\loigiai{
		Áp dụng công thức $V = \dfrac{1}{3} S_{\text{đáy}} \cdot h$, ta có: \\
$V = \dfrac{1}{3} \cdot 36 \cdot 10 = 120\text{ cm}^3$. \\
Vậy thể tích của hình chóp tam giác đều là $120\text{ cm}^3$.
	}
\end{bt}

% nguồn: De_AI_Toan8_cuoi_hoc_ki_i_1 (nhập 21/11/2026)
\begin{bt}%[8H0N38-2]%[ma:b39bc4c2]
	Cho hình chóp tam giác đều có cạnh đáy bằng $6\text{ cm}$, độ dài trung đoạn bằng $7\text{ cm}$ và diện tích mặt đáy xấp xỉ $15{,}6\text{ cm}^2$.
a) Tính diện tích xung quanh của hình chóp.
b) Tính diện tích toàn phần của hình chóp.
	\loigiai{
		a) Nửa chu vi đáy của hình chóp là $p = \dfrac{3 \cdot 6}{2} = 9\text{ cm}$. \\
Diện tích xung quanh của hình chóp là $S_{xq} = p \cdot d = 9 \cdot 7 = 63\text{ cm}^2$. \\
b) Diện tích toàn phần của hình chóp là $S_{tp} = S_{xq} + S_{\text{đáy}} \approx 63 + 15{,}6 = 78{,}6\text{ cm}^2$.
	}
\end{bt}

% nguồn: De_AI_Toan8_cuoi_hoc_ki_i_1 (nhập 21/11/2026)
\begin{bt}%[8H0N38-2]%[ma:07a9897e]
	Cho hình chóp tứ giác đều có độ dài cạnh đáy bằng $6\text{ cm}$ và độ dài trung đoạn bằng $5\text{ cm}$.
a) Tính diện tích xung quanh của hình chóp.
b) Tính diện tích mặt đáy và diện tích toàn phần của hình chóp.
	\loigiai{
		a) Nửa chu vi đáy hình vuông là $p = \dfrac{4 \cdot 6}{2} = 12\text{ cm}$. \\
Diện tích xung quanh của hình chóp là $S_{xq} = p \cdot d = 12 \cdot 5 = 60\text{ cm}^2$. \\
b) Diện tích mặt đáy là $S_{\text{đáy}} = 6^2 = 36\text{ cm}^2$. \\
Diện tích toàn phần của hình chóp là $S_{tp} = S_{xq} + S_{\text{đáy}} = 60 + 36 = 96\text{ cm}^2$.
	}
\end{bt}

% nguồn: De_AI_Toan8_cuoi_hoc_ki_i_1 (nhập 21/11/2026)
\begin{bt}%[8H0N38-2]%[ma:beeeba74]
	Cho hình chóp tứ giác đều có độ dài cạnh đáy bằng $6\text{ cm}$ và chiều cao bằng $8\text{ cm}$.
a) Tính diện tích mặt đáy của hình chóp.
b) Tính thể tích của hình chóp tứ giác đều đó.
	\loigiai{
		a) Vì mặt đáy là hình vuông cạnh $6\text{ cm}$ nên $S_{\text{đáy}} = 6^2 = 36\text{ cm}^2$. \\
b) Thể tích của hình chóp tứ giác đều là: \\
$V = \dfrac{1}{3} S_{\text{đáy}} \cdot h = \dfrac{1}{3} \cdot 36 \cdot 8 = 96\text{ cm}^3$.
	}
\end{bt}

% nguồn: De_AI_Toan8_cuoi_hoc_ki_i_1 (nhập 21/11/2026)
\begin{bt}%[8H0N38-2]%[ma:4c8af84a]
	Một hình chóp tam giác đều có thể tích bằng $80\text{ cm}^3$ và diện tích mặt đáy bằng $20\text{ cm}^2$. Tính chiều cao của hình chóp đó.
	\loigiai{
		Từ công thức thể tích hình chóp $V = \dfrac{1}{3} S_{\text{đáy}} \cdot h$, ta suy ra: \\
$h = \dfrac{3V}{S_{\text{đáy}}} = \dfrac{3 \cdot 80}{20} = 12\text{ cm}$. \\
Vậy chiều cao của hình chóp tam giác đều là $12\text{ cm}$.
	}
\end{bt}

% nguồn: De_AI_Toan8_cuoi_hoc_ki_i_1 (nhập 21/11/2026)
\begin{bt}%[8H0N38-2]%[ma:a760fb79]
	Một hình chóp tam giác đều có thể tích bằng $105\text{ cm}^3$ và chiều cao bằng $7\text{ cm}$. Tính diện tích mặt đáy của hình chóp đó.
	\loigiai{
		Từ công thức $V = \dfrac{1}{3} S_{\text{đáy}} \cdot h$, ta suy ra: \\
$S_{\text{đáy}} = \dfrac{3V}{h} = \dfrac{3 \cdot 105}{7} = 45\text{ cm}^2$. \\
Vậy diện tích mặt đáy của hình chóp là $45\text{ cm}^2$.
	}
\end{bt}
